\documentclass[10pt,a4paper]{amsart}
\usepackage[margin=19mm]{geometry}
\usepackage{amsmath,amssymb,mathtools}
\usepackage[scr=rsfs,cal=euler]{mathalfa}
\usepackage{xfrac}
\usepackage[unicode,hidelinks,bookmarks=false]{hyperref}
\newcommand{\ie}{i.e.\ }
\newcommand{\cf}{cf.\ }
\newcommand{\resp}{respectively\ }
\newcommand{\Subsection}{\subsection}
\providecommand{\nobreakdash}{\nobreak\hbox{-}}
\setlength{\emergencystretch}{2em}
\multlinegap0pt
\usepackage{mathtools}
\usepackage{amssymb}
\usepackage{tikz}
\usepackage{enumerate}
\newtheorem{proposition}{Proposition}
\newtheorem{lemma}{Lemma}
\newtheorem{corollary}{Corollary}
\def\E{\mathbb{E}}
\def \R{\mathbb {R}}
\def \T{\mathcal{T}}
\def \cL{\mathcal{L}}
\def\vp{\varepsilon}
\title{Depth profile of depth-weighted trees with bounded weights}
\author{Emmanuel Kammerer}
\date{Source: arXiv:2607.05366v1 (CC BY 4.0)}
\begin{document}
\maketitle

\noindent\emph{Source and licence.} Depth profile of depth-weighted trees with bounded weights. Version arXiv:2607.05366v1 (CC BY 4.0), distributed by arXiv under the Creative Commons Attribution 4.0 International licence, http://creativecommons.org/licenses/by/4.0/, as recorded in the arXiv licence metadata for this exact version and shown on the abstract page https://arxiv.org/abs/2607.05366v1. Full licence notice: \texttt{licenses/arxiv-2607.05366-CC-BY-4.0.txt}.\par\smallskip

\noindent\textit{Editorial scope.} Counted unit: the article's lemma nulle (source0.tex lines 561--563). The article's complete original proof of it is delivered with it (source0.tex lines 564--630, one \texttt{proof} environment of 67 lines). No mathematics is written, repaired or re-derived: the statement and the proof are the article's own lines, in the article's own notation, and no retained line is substituted or re-flowed.  Retained uncounted as the source-local context the proof needs: lines 211--213; lines 465--474; lines 525--527; lines 543--545.  Nothing was omitted from any retained span, so the package carries no omission record.  No retained line contains a citation, so the wrapper carries no bibliography and no bibliography entry.  No retained line contains a figure environment, an image-inclusion command or drawing code, so no image is delivered and the package declares no asset dependency.  Editorial formatting only, in addition: an AMS \texttt{amsart} preamble that restates the article's own declarations and macros used by the retained lines, namely the environments \texttt{proposition}, \texttt{lemma}, \texttt{corollary} on separate counters, together with the standard packages those lines need.  The retained environments print in this document as Proposition 1, Lemma 1, Corollary 1 in order of appearance, whereas the article prints the same items with the numbering of its own full document.  The complete native source of the article as distributed in the arXiv e-print for this exact version is bundled unchanged as \texttt{sources/204539/source0.tex}.
\begin{equation}\label{condition ergodicite}
	\frac{S_n}{n} \mathop{\longrightarrow}_{n\to \infty} \ell \in \R.
\end{equation}

\begin{proposition}\label{prop cv unif martingale}
	There exists $\beta_0>0$ such that the following holds for all $\beta\in (0, \beta_0)$. The sequence of functions $(z\mapsto M_n(z))_{n\ge n_0}$ converges uniformly almost surely and in $\mathrm{L}^1$ towards a random analytic function $z \mapsto M_\infty(z)$ on every compact subset of $\mathcal{E}_\beta$. Moreover, for any compact subset $K \subset \mathcal{E}_\beta$, there exists $\vp>0$ such that almost surely,
	\[
	\limsup_{n\to \infty} \max_{z \in K} n^\vp \vert M_n(z)-M_\infty(z)\vert <\infty.
	\]
	Furthermore, for any compact subset $K \subset \mathcal{E}_\beta$, for all $p>1$ small enough,
	\begin{equation}\label{eq M n borne dans L p}
	\sup_{n\ge n_0} \sup_{z \in K} \E\left( \vert M_n (z) \vert^p \right)<\infty.
	\end{equation}
\end{proposition}

\begin{lemma}\label{lemme controle hauteur totale}
	For all $\vp \in (0,1)$, the sequence of random variables $e^{(1-\vp)d(\T_n)}/n^{e/c}$ is bounded in $\mathrm{L}^1$. In particular, almost surely, for all $n$ large enough, $d(\T_n)\le (3+2e/c)\log n$.
\end{lemma}

\begin{corollary}\label{cor controle transformee de Laplace}
	Almost surely, $ \log (\cL_n(0)) \sim \log n$ as $n\to \infty$.
\end{corollary}

\begin{lemma}\label{lemme limite non nulle}
	Almost surely, $M_\infty(0)> 0$. In particular, a.s.\@ the set of $z \in (-\infty, 1)$ such that $M_\infty(z) =0$ is at most countable and its elements are isolated.
\end{lemma}

\begin{proof}
	It suffices to prove that $M_\infty(0)> 0$ almost surely since we know by Proposition \ref{prop cv unif martingale} that the function $z \mapsto M_\infty(z)$ is analytic. For all $n\ge n_0$, let $\mathcal{A}_n$ be the event that we have $d(\T_n)\le (3+2e/c)\log n$ and $\cL_n(0) \ge n^{3/4}$. For all $k\ge n_0$, let $\mathcal{B}_k \coloneqq \bigcap_{n\ge k} \mathcal{A}_n$. Note that by Lemma \ref{lemme controle hauteur totale} and Corollary \ref{cor controle transformee de Laplace}, we know that $\bigcup_{k\ge n_0}\mathcal{B}_k$ has full probability.
	
	Let $k\ge n_0$. Let us show that the sequence $(u_n)_{n\ge k}$ defined by
	\[
	u_n\coloneqq \E\left(\vert \log M_n(0)\vert^2 {\bf 1}_{\bigcap_{i=k}^n \mathcal{A}_i}\right)
	\]
	is bounded. Note that $M_n(0)>0$ for all $n\ge n_0$ so that the above display is well-defined.
	
	Let us first compute
	\begin{align*}
		\E&\left( \left. (\log M_{n+1}(0))^2{\bf 1}_{\bigcap_{i=k}^{n+1} \mathcal{A}_i} \right\vert \mathcal{F}_n \right)
		\le \E\left( \left. (\log M_{n+1}(0))^2{\bf 1}_{\bigcap_{i=k}^{n} \mathcal{A}_i} \right\vert \mathcal{F}_n \right) \\
		=&{\bf 1}_{\bigcap_{i=k}^{n} \mathcal{A}_i}
		\sum_{v \in \T_n} \frac{f(d(v))}{Z_n} \left( \log\left(\cL_n(0)+ e^{-S_{d(v)+1}}\right) -\log C_{n+1}(0)\right)^2\\
		=&{\bf 1}_{\bigcap_{i=k}^{n} \mathcal{A}_i}\sum_{v \in \T_n} \frac{f(d(v))}{Z_n} \left(\log(M_n(0)) + \log \left(1+ \frac{e^{-S_{d(v)+1}} }{\cL_n(0)}\right) -\log\left(1+\frac{1}{Z_n}\right)\right)^2\\
		=&{\bf 1}_{\bigcap_{i=k}^{n} \mathcal{A}_i} \left( (\log M_n(0))^2 + \sum_{v \in \T_n} \frac{f(d(v))}{Z_n} \left(\log \left(1+ \frac{e^{-S_{d(v)+1}} }{\cL_n(0)}\right) -\log\left(1+\frac{1}{Z_n}\right)\right)^2 \right. \\
		&\left.+ \sum_{v \in \T_n} \frac{f(d(v))}{Z_n}2(\log M_n(0))\left(\log \left(1+ \frac{e^{-S_{d(v)+1}} }{\cL_n(0)}\right) -\log\left(1+\frac{1}{Z_n}\right)\right)  \right).
	\end{align*}
	Moreover, on the event $\mathcal{A}_n$, the first sum in the last equality can be bounded from above as follows. For all $\vp>0$, there is a constant $C_\vp>0$ such that for all $n\ge k$, on the event $\mathcal{A}_n$,
	\begin{align}
		\sum_{v \in \T_n} \frac{f(d(v))}{Z_n} \left(\log \left(1+ \frac{e^{-S_{d(v)+1}} }{\cL_n(0)}\right) -\log\left(1+\frac{1}{Z_n}\right)\right)^2&\le
		\sum_{v \in \T_n} \frac{f(d(v))}{Z_n}C_\vp \left( \frac{e^{2\vp (d(\T_n)+1)}}{n^{3/2}} + \frac{1}{n^2} \right) \notag \\
		&= C_\vp \left( \frac{e^{2\vp (1+(3+2e/c)\log n)}}{n^{3/2}} + \frac{1}{n^2} \right), \label{eq premiere somme}
	\end{align}
	where in the inequality, we used \eqref{condition ergodicite} to control $S_{d(v)+1}$ and we used that $Z_n \ge cn$ since $f\ge c$.
	
	For the last sum, one can note that
	\begin{multline}\label{eq derniere somme}
		 \left\vert\sum_{v \in \T_n} \frac{f(d(v))}{Z_n}2(\log M_n(0))\left(\log \left(1+ \frac{e^{-S_{d(v)+1}} }{\cL_n(0)}\right) -\log\left(1+\frac{1}{Z_n}\right)\right) \right\vert \\ 
		 \le (1+ (\log M_n(0))^2 ) \left\vert \sum_{v \in \T_n} \frac{f(d(v))}{Z_n}\left(  \log \left(1+ \frac{e^{-S_{d(v)+1}} }{\cL_n(0)}\right) -\log\left(1+\frac{1}{Z_n}\right)\right) \right\vert 
	\end{multline}
	On the event $\mathcal{A}_n$ the sum in the absolute value on the right-hand side of \eqref{eq derniere somme} can be bounded from above by
	\[
	\sum_{v \in \T_n} \frac{f(d(v))}{Z_n} \left( \frac{e^{-S_{d(v)+1}}}{\cL_n(0)} - \log \left(1+\frac{1}{Z_n}\right)\right)= \frac{1}{Z_n} - \log \left(1+ \frac{1}{Z_n}\right) \le \frac{C}{n^2},
	\]
	for some constant $C>0$ (using that $Z_n \ge cn$ once more). Moreover, using that $\log(1+x) \ge x-x^2/2$ for all $x\ge 0$, for all $\vp>0$, there is a constant $\widetilde{C}_\vp>0$ such that for all $n\ge k$, on the event $\mathcal{A}_n$, the sum in the absolute value on the right-hand side of \eqref{eq derniere somme} can be bounded from below by
	\begin{align*}
	\sum_{v \in \T_n} &\frac{f(d(v))}{Z_n}\left( \frac{e^{-S_{d(v)+1}}}{\cL_n(0)} - \left(\frac{e^{-S_{d(v)+1}}}{\cL_n(0)} \right)^2 \right) - \log \left(1+\frac{1}{Z_n}\right)\\
	&\ge  \frac{1}{Z_n} - \log \left(1+ \frac{1}{Z_n}\right)  - \sum_{v \in \T_n} \frac{f(d(v))}{Z_n} \widetilde{C}_\vp \frac{e^{2\vp (d(\T_n)+1)}}{n^{3/2}} \\
	&\ge - \widetilde{C}_\vp \frac{e^{2\vp(1+(3+2e/c)\log n)}}{n^{3/2}}.
	\end{align*}
	Thus, for all $\vp>0$, there exists $C'_\vp>0$ such that for all $n\ge k$, on the event $\mathcal{A}_n$, 
	\begin{multline}\label{eq derniere somme 2}
		\left\vert\sum_{v \in \T_n} \frac{f(d(v))}{Z_n}2(\log M_n(0))\left(\log \left(1+ \frac{e^{-S_{d(v)+1}} }{\cL_n(0)}\right) -\log\left(1+\frac{1}{Z_n}\right)\right) \right\vert \\
		\le C'_\vp (1+ (\log M_n(0))^2) \left( \frac{e^{2\vp(1+(3+2e/c)\log n)}}{n^{3/2}} +  \frac{1}{n^2} \right).
	\end{multline}
	Combining \eqref{eq derniere somme 2} with \eqref{eq premiere somme}, and taking $\vp>0$ small enough, we deduce that there exists a constant $K>0$ such that for all $n\ge k$,
	\[
	\E\left(\left. (\log M_{n+1}(0))^2 {\bf 1}_{\bigcap_{i=k}^{n+1} \mathcal{A}_i} \right\vert \mathcal{F}_n\right)\le {\bf 1}_{\bigcap_{i=k}^n \mathcal{A}_i} \left( \left(1+ \frac{K}{n^{5/4}} \right) (\log M_n(0))^2 + \frac{K}{n^{5/4}} \right).
	\]
	Taking the expectation, we deduce that for all $n\ge k$,
	\[
	u_{n+1} \le \left(1 + \frac{K}{n^{5/4}}\right) u_n + \frac{K}{n^{5/4}}.
	\]
	One can then easily see that $(u_n)_{n\ge k}$ is bounded (letting $v_n = 1+u_n$, we see that $v_{n+1} \le (1+K/n^{5/4})v_n$).
	
	Recall that $\mathcal{B}_k = \bigcap_{n\ge k} \mathcal{A}_n$. Note that for all $n\ge k$,
	\[
	\E\left((\log M_n(0))^2 {\bf 1}_{\mathcal{B}_k} \right)\le u_n.
	\]
	In particular, by Fatou's lemma, since $M_n(0)\to M_\infty(0)$ almost surely by Proposition \ref{prop cv unif martingale}, we conclude that
	\[
	\E\left((\log M_\infty(0))^2  {\bf 1}_{\mathcal{B}_k} \right) \le \liminf_{n\to \infty}\E\left((\log M_n(0))^2 {\bf 1}_{\mathcal{B}_k} \right) <\infty. 
	\]
	Thus, almost surely, $M_\infty(0) >0$ on the event $\mathcal{B}_k$. Letting $k\to \infty$ proves that $M_\infty(0)>0$ almost surely. Finally, recall from Proposition \ref{prop cv unif martingale} that $z\mapsto M_\infty(z)$ is a random analytic function (as a uniform limit of analytic functions). Therefore, since the zeros of a non-zero analytic function are at most countable and isolated, we conclude that almost surely, on the interval $(-\infty, 1)$, the function $z\mapsto M_\infty(z)$ has at most countably many zeros that are isolated.
\end{proof}

\end{document}
